\documentclass[10pt,a4paper]{amsart}
\usepackage{amsmath,amssymb,amsthm,graphicx}
\usepackage{verbatim}
\newenvironment{reference}{\comment}{\endcomment}
%\newenvironment{reference}{}{}
\newenvironment{slogan}{\comment}{\endcomment}
\newenvironment{history}{\comment}{\endcomment}

% For commutative diagrams we use Xy-pic
\usepackage[all]{xy}

% We use 2cell for 2-commutative diagrams.
\xyoption{2cell}
% Theorem environments.
%
\theoremstyle{plain}
\newtheorem{theorem}[subsection]{Theorem}
\newtheorem{proposition}[subsection]{Proposition}
\newtheorem{lemma}[subsection]{Lemma}

\theoremstyle{definition}
\newtheorem{definition}[subsection]{Definition}
\newtheorem{example}[subsection]{Example}
\newtheorem{exercise}[subsection]{Exercise}
\newtheorem{situation}[subsection]{Situation}

\theoremstyle{remark}
\newtheorem{remark}[subsection]{Remark}
\newtheorem{remarks}[subsection]{Remarks}

\numberwithin{equation}{subsection}

% Macros
%
\def\lim{\mathop{\mathrm{lim}}\nolimits}
\def\colim{\mathop{\mathrm{colim}}\nolimits}
\def\Spec{\mathop{\mathrm{Spec}}}
\def\Hom{\mathop{\mathrm{Hom}}\nolimits}
\def\Ext{\mathop{\mathrm{Ext}}\nolimits}
\def\SheafHom{\mathop{\mathcal{H}\!\mathit{om}}\nolimits}
\def\SheafExt{\mathop{\mathcal{E}\!\mathit{xt}}\nolimits}
\def\Sch{\mathit{Sch}}
\def\Mor{\mathop{\mathrm{Mor}}\nolimits}
\def\Ob{\mathop{\mathrm{Ob}}\nolimits}
\def\Sh{\mathop{\mathit{Sh}}\nolimits}
\def\NL{\mathop{N\!L}\nolimits}
\def\CH{\mathop{\mathrm{CH}}\nolimits}
\def\proetale{{pro\text{-}\acute{e}tale}}
\def\etale{{\acute{e}tale}}
\def\QCoh{\mathit{QCoh}}
\def\Ker{\mathop{\mathrm{Ker}}}
\def\Im{\mathop{\mathrm{Im}}}
\def\Coker{\mathop{\mathrm{Coker}}}
\def\Coim{\mathop{\mathrm{Coim}}}

% Boxtimes
%
\DeclareMathSymbol{\boxtimes}{\mathbin}{AMSa}{"02}

%
% Macros for moduli stacks/spaces
%
\def\QCohstack{\mathcal{QC}\!\mathit{oh}}
\def\Cohstack{\mathcal{C}\!\mathit{oh}}
\def\Spacesstack{\mathcal{S}\!\mathit{paces}}
\def\Quotfunctor{\mathrm{Quot}}
\def\Hilbfunctor{\mathrm{Hilb}}
\def\Curvesstack{\mathcal{C}\!\mathit{urves}}
\def\Polarizedstack{\mathcal{P}\!\mathit{olarized}}
\def\Complexesstack{\mathcal{C}\!\mathit{omplexes}}
% \Pic is the operator that assigns to X its picard group, usage \Pic(X)
% \Picardstack_{X/B} denotes the Picard stack of X over B
% \Picardfunctor_{X/B} denotes the Picard functor of X over B
\def\Pic{\mathop{\mathrm{Pic}}\nolimits}
\def\Picardstack{\mathcal{P}\!\mathit{ic}}
\def\Picardfunctor{\mathrm{Pic}}
\def\Deformationcategory{\mathcal{D}\!\mathit{ef}}
\usepackage{cleveref}
\crefname{originalreference}{original reference}{original references}
\Crefname{originalreference}{Original reference}{Original references}
\makeatletter
\newlabel{stacks-sheaves-section-quasi-coherent-algebraic-stacks}{{06WU}{source}{}{sourcepointer-1}{}}
\newlabel{stacks-sheaves-section-quasi-coherent-algebraic-stacks@cref}{{[originalreference][0][]06WU}{[1][0][]source}{}{}{}}
\makeatother
\crefname{originalreference}{original reference}{original references}
\Crefname{originalreference}{Original reference}{Original references}
\makeatletter
\newlabel{spaces-groupoids-lemma-abelian}{{06VZ}{source}{}{sourcepointer-2}{}}
\newlabel{spaces-groupoids-lemma-abelian@cref}{{[originalreference][0][]06VZ}{[1][0][]source}{}{}{}}
\makeatother
\crefname{originalreference}{original reference}{original references}
\Crefname{originalreference}{Original reference}{Original references}
\makeatletter
\newlabel{stacks-sheaves-lemma-quasi-coherent}{{06WK}{source}{}{sourcepointer-3}{}}
\newlabel{stacks-sheaves-lemma-quasi-coherent@cref}{{[originalreference][0][]06WK}{[1][0][]source}{}{}{}}
\makeatother
\title{Injectivity and surjectivity of qc maps checked on flat coverings}
\author{The Stacks Project Authors}
\date{stacks-project:stacks-cohomology.tex:local-snapshot-20261012; GNU FDL 1.2 or later}
\begin{document}
\maketitle
\noindent\textit{Editorial note.} All delivered mathematical text is the author's own, copied verbatim from the pinned Stacks Project chapter file stacks-cohomology.tex; no statement or proof step was rewritten, reordered or removed.\par\medskip
\noindent\textit{Editorial note.} The counted span cites the lemmas recorded in assembly.external\_references with canonical tags as navigation pointers.\par\medskip
\noindent\textit{Allowed standard prerequisites (editorial).}\begin{itemize}
\item Abelian structure on groupoid modules. Groupoids in Spaces, Lemma spaces-groupoids-lemma-abelian. Exactness detected on the presentation cover.
\item Pullback and restriction of qc modules. Sheaves on Stacks, Lemma stacks-sheaves-lemma-quasi-coherent; Definition stacks-sheaves-definition-pullback. Compatibility of pullbacks with restrictions.
\item Flat pullback reflects exactness. Standard scheme fact: flat pullback is exact; jointly surjective flat families reflect injective/surjective maps of modules (checked at stalks). Scheme-level reflection step.
\end{itemize}\medskip
\ifdefined\hypertarget\hypertarget{sourcepointer-1}{}\fi
\noindent\textit{Source reference.} Original reference 06WU: sources/stacks-chow/tags.txt, source lines 8828--8828. Cross-chapter standard prerequisite cited by the author inside the counted span; it is used as a pointer and its own proof is not part of this item.\par\medskip
\ifdefined\hypertarget\hypertarget{sourcepointer-2}{}\fi
\noindent\textit{Source reference.} Original reference 06VZ: sources/stacks-chow/tags.txt, source lines 8798--8798. Cross-chapter standard prerequisite cited by the author inside the counted span; it is used as a pointer and its own proof is not part of this item.\par\medskip
\ifdefined\hypertarget\hypertarget{sourcepointer-3}{}\fi
\noindent\textit{Source reference.} Original reference 06WK: sources/stacks-chow/tags.txt, source lines 8819--8819. Cross-chapter standard prerequisite cited by the author inside the counted span; it is used as a pointer and its own proof is not part of this item.\par\medskip
\begin{lemma}
\label{lemma-quasi-coherent-check-exact}
Let $\mathcal{X}$ be an algebraic stack. Let $I$ be a set and
for $i \in I$ let $x_i : U_i \to \mathcal{X}$ be an object
of $\mathcal{X}$. Assume that $x_i$ is flat and
$\coprod x_i : \coprod U_i \to \mathcal{X}$ is surjective.
Let $\varphi : \mathcal{F} \to \mathcal{G}$ be an arrow of
$\QCoh(\mathcal{O}_\mathcal{X})$. Denote $\varphi_i$
the restriction of $\varphi$ to $(U_i)_\etale$.
Then $\varphi$ is injective, resp.\ surjective, resp.\ an isomorphism
if and only if each $\varphi_i$ is so.
\end{lemma}

\begin{proof}
Choose a scheme $U$ and a surjective smooth morphism
$x : U \to \mathcal{X}$. We may and do think of $x$
as an object of $\mathcal{X}$.
This produces a presentation $\mathcal{X} = [U/R]$ for some groupoid
in spaces $(U, R, s, t, c)$ and correspondingly an equivalence
$$
\QCoh(\mathcal{O}_\mathcal{X}) = \QCoh(U, R, s, t, c)
$$
See discussion in Sheaves on Stacks, Section
\ref{stacks-sheaves-section-quasi-coherent-algebraic-stacks}.
The structure of abelian category on the right hand is such that $\varphi$
is injective, resp.\ surjective, resp.\ an isomorphism
if and only if the restriction $\varphi|_{U_\etale}$ is so, see
Groupoids in Spaces, Lemma \ref{spaces-groupoids-lemma-abelian}.

\medskip\noindent
For each $i$ we choose an \'etale covering
$\{W_{i, j} \to V \times_\mathcal{X} U_i\}_{j \in J_i}$
by schemes. Denote $g_{i, j} : W_{i, j} \to V$ and
$h_{i, j} : W_{i, j} \to U_i$ the obvious arrows.
Each of the morphisms of schemes $g_{i, j} : W_{i, j} \to U$ is flat
and they are jointly surjective.
Similarly, for each fixed $i$ the morphisms of schemes
$h_{i, j} : W_{i, j} \to U_i$ are flat and jointly surjective.
By Sheaves on Stacks, Lemma \ref{stacks-sheaves-lemma-quasi-coherent}
the pullback by $(g_{i, j})_{small}$ of the restriction
$\varphi|_{U_\etale}$ is the restriction $\varphi|_{(W_{i, j})_\etale}$
and the pullback by $(h_{i, j})_{small}$ of the restriction
$\varphi|_{(U_i)_\etale}$ is the restriction $\varphi|_{(W_{i, j})_\etale}$.
Pullback of quasi-coherent modules by a flat morphism of schemes is exact
and pullback by a jointly surjective family of flat morphisms of schemes
reflects injective, resp.\ surjective, resp.\ bijective maps
of quasi-coherent modules (in fact this holds for all modules
as we can check exactness at stalks). Thus we see
$$
\varphi|_{U_\etale} \text{ injective}
\Leftrightarrow
\varphi|_{(W_{i, j})_\etale} \text{ injective for all }i, j
\Leftrightarrow
\varphi|_{(U_i)_\etale} \text{ injective for all }i
$$
This finishes the proof.
\end{proof}

\end{document}
